\section{数列与有限求和}
\label{sec:ru-sequences-finite-sums}

数列是以自然数为输入的函数。把它单独列出，是因为下一章的极限首先在数列上定义，而无限求和又要从有限和出发。
本节说明数列与集合、函数族的区别，再推导几类不需要极限的有限和公式。

一个实数\term{数列}（sequence）就是映射$a:\{1,2,\ldots\}\to\mathbb R$，把$n$处的值写作$a_n$，整个数列记作$(a_n)_{n\geq1}$或$(a_n)$；有时下标从$0$开始。
数列记录了每一项的位置，因而与它的取值集合不同：数列$((-1)^n)$在$-1$与$1$之间无限次交替，取值集合$\{-1,1\}$却只有两个元素，次序和重复都看不出来。
数列也不同于函数族：数列的每一项是一个数，函数族的每个元素是一整个函数。下一章研究函数列时，每个编号对应的才是一个函数。

第~\ref{sec:prelim-finite-sums-products}节已经约定了有限和$\sum_{k=1}^na_k$与有限积的记号，这里只回顾三条后面常用的规则。
平移哑索引时范围随之改写，如$\sum_{k=1}^{n}a_{k+1}=\sum_{j=2}^{n+1}a_j$；空和为$0$、空积为$1$，使拆分求和时不必单独处理空的部分；
有限个项的二重和可以交换次序，内层范围依赖外层索引时按式~\eqref{eq:prelim-triangular-sum}重写范围。
交换次序的依据是每个索引对恰好被计入一次，这只对有限多项成立；无限多项的和要先用极限定义，第~\ref{sec:ru-series-rearrangement-grouping}节将看到交换次序可能改变结果。

几类有限和可以直接写出公式。先看$\sum_{k=1}^nk$：把它正着写一遍、倒着写一遍再逐项相加，每一对都等于$n+1$，共$n$对，所以
\[
 \sum_{k=1}^nk=\frac{n(n+1)}2.
\]
由此得到首项为$a$、公差为$d$的等差数列前$n$项之和：
\[
 \sum_{k=0}^{n-1}(a+kd)=na+d\sum_{k=0}^{n-1}k=na+\frac{n(n-1)}2d.
\]
公比为$\rho$的几何和由式~\eqref{eq:prelim-finite-geometric-sum}给出：$\rho\neq1$时$\sum_{k=0}^N\rho^k=(1-\rho^{N+1})/(1-\rho)$，$\rho=1$时为$N+1$。
两个公式都要写清首项、项数和公比；求和从$k=1$开始时，几何和要再乘一个$\rho$。

若每一项能写成相邻两项之差$a_k=b_k-b_{k+1}$，式~\eqref{eq:prelim-telescoping-sum}给出$\sum_{k=1}^Na_k=b_1-b_{N+1}$，中间各项两两抵消，称为裂项求和。
例如$\frac1{k(k+1)}=\frac1k-\frac1{k+1}$，所以
\[
 \sum_{k=1}^N\frac1{k(k+1)}=1-\frac1{N+1}.
\]
这些恒等式对每个固定的$N$成立，只用到有限次运算。当$N$无限增大时，$1-1/(N+1)$与$(1-\rho^{N+1})/(1-\rho)$会怎样变化，下一章用极限回答。
